% 全笔记统一数学宏
% 设计原则：
% 1. 用简短命令减少重复输入和排版差异；
% 2. 所有宏使用 \providecommand，避免覆盖已有定义；
% 3. 在可能产生歧义的地方保留显式写法（如 \mathbb{R}^d）。

% ---- 常用数集 ----
\providecommand{\R}{\mathbb{R}}
\providecommand{\N}{\mathbb{N}}
\providecommand{\Z}{\mathbb{Z}}
\providecommand{\Q}{\mathbb{Q}}
\providecommand{\C}{\mathbb{C}}

% ---- 概率与期望 ----
\providecommand{\E}{\mathbb{E}}
% 概率测度全书统一用斜体 P（见 STYLE_GUIDE.md 第 6 节）；\Prob 保留为兼容别名
\providecommand{\Prob}{P}
\providecommand{\Var}{\operatorname{Var}}
\providecommand{\Cov}{\operatorname{Cov}}
\providecommand{\Corr}{\operatorname{Corr}}
\providecommand{\Bias}{\operatorname{Bias}}
\providecommand{\SE}{\operatorname{SE}}
\providecommand{\MSE}{\operatorname{MSE}}
\providecommand{\ind}{\mathbf{1}}          % 指示函数

% ---- 优化算子 ----
\providecommand{\argmax}{\operatorname*{arg\,max}}
\providecommand{\argmin}{\operatorname*{arg\,min}}

% ---- 线性代数 ----
\providecommand{\trans}{\mathsf{T}}        % 转置
\providecommand{\conj}{\mathsf{H}}         % 共轭转置
\providecommand{\trace}{\operatorname{tr}}
\providecommand{\diag}{\operatorname{diag}}
\providecommand{\rank}{\operatorname{rank}}
\providecommand{\nullsp}{\operatorname{Null}}
\providecommand{\rangesp}{\operatorname{Range}}

% ---- 范数、绝对值、内积、集合 ----
% 这些宏自动加 \left...\right，但在行内小公式中可能偏宽。
% 如果追求紧凑，可以改用 \norm*、\abs*（需要 mathtools 的 \DeclarePairedDelimiter）。
\providecommand{\norm}[1]{\left\lVert#1\right\rVert}
\providecommand{\abs}[1]{\left|#1\right|}
\providecommand{\ip}[2]{\left\langle#1,#2\right\rangle} % inner product
\providecommand{\set}[1]{\left\{#1\right\}}
\providecommand{\given}{\mid}

% ---- 常用算子 ----
\providecommand{\sign}{\operatorname{sign}}
\providecommand{\relu}{\operatorname{ReLU}}
\providecommand{\softmax}{\operatorname{softmax}}
\providecommand{\logsumexp}{\operatorname{LSE}}
\providecommand{\KL}{D_{\mathrm{KL}}}
\providecommand{\TV}{\operatorname{TV}}

% ---- 标注与强调 ----
\providecommand{\eqdef}{\stackrel{\text{\tiny def}}{=}} % 按定义相等
